\documentclass[11pt]{report}
\usepackage[english]{babel} 
\usepackage{mathpazo}%[utf8]{inputenc}

%%%PAGE DIMENSIONS
%\usepackage{geometry} % to change the page dimensions
%\geometry{a4paper} % or letterpaper (US) or a5paper or....
%%%IN ALTO ALLA PAGINA...

%%% PER FRONTESPIZIO, poi nel corpo della tesi cambiamo impaginazione...
\usepackage{graphicx}
\usepackage[a4paper]{geometry} % per modificare le dimensioni e i margini del documento	%NB:  DA  FRONTESPIZIO DA UNI
\newgeometry{top=2cm, bottom=1.8cm, left=2cm, right=2cm} 						%NB:  PER FRONTESPIZIO DA UNI
%%%%%%%%%%%%%%%%%%%%%%%%%%%%%%%%%%%%%%

\graphicspath{{images/}}
%%%Bibliografia  ---->  \cite{...}, \parencite{...} o \textcite{...}
\usepackage[sorting=none]{biblatex}
%\usepackage[backend=biber, style=authoryear, sorting=nyt]{biblatex} % sicché Stile Autore-Anno classico (es. Rossi 2020)
\addbibresource{References_fromZotero.bib} % Nome del tuo file .bib con estensione
\usepackage{hyperref}
%%%

\usepackage{lipsum}
%%% MIE AGGIUNTE
\usepackage{amsthm} % to have proof automatically
\newtheorem{remark}{Remark}
%\newtheorem{theorem}{Proposizione}
%\newtheorem{definition}{Definizione}

\usepackage{amsmath}  %  equazioni molto lunghe  e includere molte equazioni nella stessa linea
\usepackage{amssymb}
\usepackage{empheq}
\usepackage{gensymb} % per usare "\degree" "\ohm" "\perthousand"
\usepackage{siunitx } % notazione scientifica "\num{3.12e-2}"
\usepackage{cancel}
\usepackage{xcolor}

\usepackage{titlesec}

% Rimuove la scritta "Capitolo" e mostra solo il numero seguito dal titolo
\titleformat{\chapter}[hang]
  {\normalfont\huge\bfseries}
  {\thechapter\quad} % Mostra solo il numero (es. "1   Titolo")
  {0pt}
  {\Huge}

%%% END Article customizations

%%% The "real" document content comes below...

%\title{
%		{\Huge Metodi matematici per la Fisica}\\ {\large (9 CFU)}\\
%		{\large appunti dalle lezioni del prof F. Benatti} \\
%		%{\includegraphics{university.jpg}}
%}
%\author{Giuseppe Cipolla}
%\date{a.a. 2017-2018}

\begin{document}
%\maketitle
% !TEX encoding = UTF-8
% !TEX program = pdflatex
% !TEX root = ../A-thesis-main.tex

%\documentclass{article}

%\usepackage{graphicx} % per inserire immagini
%\usepackage[a4paper]{geometry} % per modificare le dimensioni e i margini del documento



%\begin{document}

% MODIFICARE PER CAMBIARE I MARGINI
%\newgeometry{top=2cm, bottom=1.8cm, left=2cm, right=2cm}

%\fontfamily{pbk}\selectfont

\thispagestyle{empty} 

\begin{center}

\begin{figure}
    \centering
    \includegraphics[height=5.44cm]{Sigillo_Scritta_HR.png}
\end{figure}


\hrule

\vspace{1cm}

\LARGE{\textbf{Dipartimento di Fisica}}

\vspace{0.5cm}

\LARGE{\textbf{CORSO DI LAUREA MAGISTRALE IN FISICA}}

\vspace{0.5cm}

\LARGE{\textbf{\emph{Curriculum Fisica  Teorica}}}
\vspace{0.5cm}

\LARGE{\textbf{TESI DI LAUREA MAGISTRALE}}

\vspace{.5cm}

\vspace{1.0cm}
\huge{\textbf{Duis aute irure reprehenderit\\in voluptate velit esse cillum\\dolore eu fugiat nulla \emph{pariatur}}}

\end{center}

\vfill

\begin{minipage}[t]{0.43\textwidth}
\begin{flushleft}


\large{\emph{Laureando:}}

\large{\textbf {Giuseppe Cipolla}} 

\end{flushleft}
\end{minipage}
\begin{minipage}[t]{0.5\textwidth}
\begin{flushright}
\large{\emph{Relatore:}}

\large{\textbf {Prof. Nicolas Brunel}} 

\vspace{.5cm}

%\large{\emph{Correlatori:}}

%\large{\textbf{Prof. Roberto Rossi\\Prof. Giacomo Grigio}}

\end{flushright}
\end{minipage}

\vspace{2cm}

\begin{center}
    \emph{\textbf{ANNO ACCADEMICO 2026-2027}}
\end{center}

\restoregeometry

%\end{document}

% Inserisci \newgeometry subito dopo il frontespizio
\newgeometry{
    top=3cm, 
    bottom=2.5cm, 
    left=2.7cm,   % Margine sinistro più ampio per la rilegatura?
    right=2.5cm
}

%%\lipsum
%\chapter*{Dedication}
%to mum and dad...
% Pasquale, Miriam, Salvatore D'Angelo, prof Giugliano, prof Marsili, prof Brunel 

%\chapter*{Declaration}
%I declare that...
%\chapter*{Acknowledgement}
%I want to thank...

\chapter*{Abstract \& Sommario}
%\addcontentsline{toc}{chapter}{Abstract and Summary}

\noindent\textbf{Abstract (English)} \\[0.5em]
This thesis approaches the problem of temporal sequences disambiguation in the context of neural network -  specifically rate neuron models - within neural populations by integrating and generalizing two previosly enstablished connectivity frameworks: one leading to fixed attractor dynamics \cite{pereira2018}, whereas the other generates sequential activity \cite{gillett2020}, even with tunable retrieval speed \cite{gillett2024}. In order to characterize the disambiguation of two different temporal sequences starting from the same initial pattern of activity, we complement numerical simulations with the analytical derivation of the Dynamic Mean Field (DMF) equations for the system via a Path Integral formulation. This approach, indeed, allows to derive also the spectrum of fluctuations around the mean field equations, from which we derive the general stability condition for the DMF states. \cite{cris&somp2018, metz2025}.

\vspace{1.5cm}

\noindent\textbf{Sommario (Italiano)} \\[0.5em]
Questo elaborato di tesi si concentra sullo studio e sull'analisi di un modello di rete neurale che mette in relazione due modelli di reti neurali con connettività già studiata in letteratura, con l'obiettivo di studiare lo spazio dei parametri che permette di ottenere disambiguazione di sequence temporali. L'analisi quantitativa è basata sulla derivazioni delle equazioni di campo medio passando per il formalismo a path integral, che permette anche di ottenere lo spettro delle fluttuazioni rispetto alle equazioni di campo medio, da cui si derivano in generale le condizioni di generale stabilità degli stati di 

\tableofcontents

\chapter{Introduction}
% !TEX encoding = UTF-8
% !TEX program = pdflatex
% !TEX root = ../A-thesis-main.tex







\chapter{Theory}
Draft: intro to the model and to the DMFT, and Hopfield-like vrs of the model ?
% !TEX encoding = UTF-8
% !TEX program = pdflatex
% !TEX root = ../A-thesis-main.tex

\section{The Model}
We consider rate model.... ecc ecc, focusing for computaional convenience on the dynamics of the field $h$

\subsection{h-dynamics}
\begin{equation} \label{model_eqq_hdyn}
\begin{pmatrix} \tau_F\,\mathbb{1}_{N_F\times N_F} & \mathbb{O}_{N_F\times N_S} \\ \mathbb{O}_{N_S\times N_F} & \tau_S\,\mathbb{1}_{N_S\times N_S} \end{pmatrix} \cfrac{d}{dt}\begin{pmatrix}
\vec{h}^F \\ \vec{h}^S \end{pmatrix} = -\begin{pmatrix}
\vec{h}^F \\ \vec{h}^S \end{pmatrix}+\begin{pmatrix} J^{FF}_{N_{F}\times N_{F}} & J^{FS}_{N_{F}\times N_{S}} \\ J^{SF}_{N_{S}\times N_{F}} & J^{SS}_{N_{S}\times N_{S}}\end{pmatrix}\underbrace{\phi\begin{pmatrix}
\vec{h}^F \\ \vec{h}^S \end{pmatrix}}_{\text{element-wise}}
\end{equation}
where 
\begin{equation}
\begin{aligned}
J^{FF}_{ij}&:=\cfrac{A_{FF}}{K_{FF}}c_{ij}\sum_{s=1}^{\mathcal{S}}\limits\sum_{\mu=1}^{\mathcal{P}-1}\limits f(\xi^{s,\mu+1}_{i})g(\xi^{s,\mu}_{j}) && J^{FS}_{ik}:=\cfrac{A_{FS}}{K_{FS}}c_{ik}\sum_{s=1}^{\mathcal{S}}\limits\sum_{\mu=1}^P\limits f(\xi^{s,\mu}_{i})g(\eta^{s}_{k}) \\ J^{SF}_{ij}&:=\cfrac{A_{SF}}{K_{SF}}c_{ij}\sum_{s=1}^{\mathcal{S}}\limits\sum_{\mu=1}^{\mathcal{P}}\limits f(\eta^{s}_{i})g(\xi^{s,\mu}_{j}) && J^{SS}_{ik} :=\cfrac{A_{SS}}{K_{SS}}c_{ik}\sum_{s=1}^{\mathcal{S}}\limits f(\eta^{s}_{i})g(\eta^{s}_{k})
\end{aligned}
\end{equation}
with:
\begin{itemize}
\item $\mathcal{S}=2$ (the number of temporal sequences), and $\mathcal{P}=16$ (the number of patterns for each temporal sequence);
\item $K_{target,source}:=\mathbb{E}[N_{\textcolor{red}{source}}\,c_{ij}]$ , in red because in the code (\texttt{connectivity.py}, line 197: \texttt{self.K = p*target.size`)}\\
	\textbf{Remark}: $N_{source}=N_{target}$ for the moment...(one problem at a time!) 
\item $f$, and $g$, step functions as in Gillett 2020.
\end{itemize}
\subsection{A step behind: the Hopfield-like version of the model}
In order to get an insight on the network behavior, we consider the following simplest version of our h-dynamics. A time-discrete dynamical system where each $h_i(t)$ is now $S_i(k)\in \{\pm 1\}$, stored pattern $\xi,\eta\in\{\pm 1\}$, and the time-evolution is given by the iterative map
$ S_i(k+1)=\text{sgn}\big(\sum_{j\neq i}W_{ij}S_j(k)\big)$.  Thus, Eq \eqref{model_eqq_hdyn} for $A_{SF}=A_{FS}\equiv A$, $A_{FF},A_{SS}>0$, leads to
the following update rule for the fast population states $S_{i_F}^F(t+1)$:
\begin{equation}
S_{i_F}^F(t+1) = \operatorname{sgn}\left( \sum_{j_F=1}^{N_F} W_{i_F j_F}^{FF} S_{j_F}^F(t) + \sum_{j_S=1}^{N_S} W_{i_F j_S}^{FS} S_{j_S}^S(t) \right)
\end{equation}
where the synaptic coupling matrices are defined as:
\begin{align}
W_{i_F j_F}^{FF} &= A_{FF} \frac{1}{N_F} \sum_{s=1}^{\mathcal{S}} \sum_{\mu=1}^{\mathcal{P}-1} \xi_{i_F}^{s,\mu+1} \xi_{j_F}^{s,\mu} &
W_{i_F j_S}^{FS} &= A_{FS} \frac{1}{N_S} \sum_{s=1}^{\mathcal{S}} \sum_{\mu=1}^{\mathcal{P}} \xi_{i_F}^{s,\mu} \eta_{j_S}^{s}
\end{align}
Substituting the couplings into the update equation and rearranging the sums (towards the application of the Law of Large Numbers, LLN), we obtain:
\begin{align}
S_{i_F}^F(t+1) &= \operatorname{sgn}\left[ \frac{A_{FF}}{A} \sum_{j_F=1}^{N_F} \left( \frac{1}{N_F} \sum_{s=1}^{S} \sum_{\mu=1}^{P-1} \xi_{i_F}^{s,\mu+1} \xi_{j_F}^{s,\mu} \right) S_{j_F}^F(t) + \sum_{j_S=1}^{N_S} \left( \frac{1}{N_S} \sum_{s=1}^{S} \sum_{\mu=1}^{P} \xi_{i_F}^{s,\mu} \eta_{j_S}^{s} \right) S_{j_S}^S(t) \right] \nonumber \\
&= \operatorname{sgn}\left[ \frac{A_{FF}}{A} \sum_{s=1}^{S} \sum_{\mu=1}^{P-1} \left( \frac{1}{N_F} \sum_{j_F=1}^{N_F} \xi_{j_F}^{s,\mu} S_{j_F}^F(t) \right) \xi_{i_F}^{s,\mu+1} + \sum_{s=1}^{S} \sum_{\mu=1}^{P} \left( \frac{1}{N_S} \sum_{j_S=1}^{N_S} \eta_{j_S}^{s} S_{j_S}^S(t) \right) \xi_{i_F}^{s,\mu} \right] \tag{$\star$}
\end{align}

\subsubsection*{Initial Conditions and First Step ($t=0 \to t=1$)}

Now, suppose that the initial patterns of the two sequences coincide:
\begin{equation}
\vec{\xi}^{1,1} = \vec{\xi}^{2,1} \quad (\text{i.e. } \xi_k^{1,1} = \xi_k^{2,1}, \; \forall k=1,\dots,N_F)
\end{equation}
In addition, at $t=0$, we initialize the fast population (\textsc{fast pop}) and slow population as:
\begin{equation}
t=0: \quad
\begin{cases}
\vec{S}^F(0) = \vec{\xi}^{2,1} & (\text{or } \vec{\xi}^{1,1}, \text{ it's the same}) \\
\vec{S}^S(0) = \vec{\eta}^{x}, \quad x \in \{1,2\} & (\text{index for the temporal sequence})
\end{cases}
\end{equation}

Evaluating equation $(\star)$ at $t=0$:
\begin{equation}
(\star)\big|_{t=0}: \quad S_{i_F}^F(1) = \operatorname{sgn}\left[ \frac{A_{FF}}{A} \sum_{s=1}^{S} \sum_{\mu=1}^{P-1} \left( \frac{1}{N_F} \sum_{j_F=1}^{N_F} \xi_{j_F}^{s,\mu} \xi_{j_F}^{2,1} \right) \xi_{i_F}^{s,\mu+1} + \sum_{s=1}^{S} \sum_{\mu=1}^{P} \left( \frac{1}{N_S} \sum_{j_S=1}^{N_S} \eta_{j_S}^{s} \eta_{j_S}^{x} \right) \xi_{i_F}^{s,\mu} \right]
\end{equation}

In the thermodynamic limit ($N_F \to +\infty$, $N_S \to +\infty$), by the Law of Large Numbers (LLN):
\begin{align}
\frac{1}{N_F} \sum_{j_F=1}^{N_F} \xi_{j_F}^{s,\mu} \xi_{j_F}^{2,1} &\xrightarrow[N_F \to +\infty]{\text{LLN}} \mathbb{E}\left[\xi^{s,\mu}\xi^{2,1}\right] = \delta_{\mu,1} \quad (\text{for } s \in \{1,2\}, \text{ since } \vec{\xi}^{1,1} = \vec{\xi}^{2,1}) \\
\frac{1}{N_S} \sum_{j_S=1}^{N_S} \eta_{j_S}^{s} \eta_{j_S}^{x} &\xrightarrow[N_S \to +\infty]{\text{LLN}} \mathbb{E}\left[\eta^{s}\eta^{x}\right] = \delta_{s,x}
\end{align}

Thus, for $S=2$ sequences and $x \in \{1,2\}$, we obtain:
\begin{equation}
S_{i_F}^F(1) \xrightarrow[N_F, N_S \to +\infty]{\text{LLN}} \operatorname{sgn}\left( \frac{A_{FF}}{A} \sum_{s=1}^{S} \xi_{i_F}^{s,2} + \sum_{\mu=1}^{P} \xi_{i_F}^{x,\mu} \right)
\end{equation}
\textit{Note on moments:} For $\xi \in \{-1, 1\}$ with $\operatorname{prob}(\xi = \pm 1) = 1/2$, we have $\mu = 0$ and $\sigma^2 = 1$, which implies $\operatorname{sgn}(\xi) = \xi$ (same distribution!).

\subsubsection*{Regime Analysis}

We analyze the behavior of $S_{i_F}^F(1)$ across three regimes of the ratio $\frac{A_{FF}}{A}$:

\begin{enumerate}
    \item $\dfrac{A_{FF}}{A} \gg 1$:
    \begin{equation}
    S_{i_F}^F(1) \approx \operatorname{sgn}\left( \sum_{s=1}^{S} \xi_{i_F}^{s,2} \right)
    \end{equation}
    \textbf{Random retrieval} of $\xi_{i_F}^{1,2}$ or $\xi_{i_F}^{2,2}$, i.e. slow population does not affect fast population which can do back and forth from one sequence to the other.

    \item $\dfrac{A_{FF}}{A} \ll 1$:
    \begin{equation}
    S_{i_F}^F(1) \approx \operatorname{sgn}\left( \sum_{\mu=1}^{P} \xi_{i_F}^{x,\mu} \right)
    \end{equation}
    \textbf{Fully disambiguation} (guided by the slow population), \textbf{but loss of the temporal sequence}: equiprobability of each pattern of the selected temporal sequence.

    \item $\dfrac{A_{FF}}{A} \simeq 1$:
    \begin{align}
    S_{i_F}^F(1) &\approx \operatorname{sgn}\left( \sum_{s=1}^{S} \xi_{i_F}^{s,2} + \sum_{\mu=1}^{P} \xi_{i_F}^{x,\mu} \right) \nonumber \\
    &= \operatorname{sgn}\left( \xi_{i_F}^{1,2} + \xi_{i_F}^{2,2} + \xi_{i_F}^{x,1} + \xi_{i_F}^{x,2} + \sum_{\mu=3}^{P} \xi_{i_F}^{x,\mu} \right) \quad (\text{extracting the first two terms})
    \end{align}
    \textbf{Unifying Notation:} If $x=1$, then $\bar{x}=2$; if $x=2$, then $\bar{x}=1$ (wlog, let's suppose $x=2$). Rewriting the expression:
    \begin{align}
    S_{i_F}^F(1) &\approx \operatorname{sgn}\left( \xi_{i_F}^{\bar{x},2} + \xi_{i_F}^{x,2} + \xi_{i_F}^{x,1} + \xi_{i_F}^{x,2} + \sum_{\mu=3}^{P} \xi_{i_F}^{x,\mu} \right) \nonumber \\
    &= \operatorname{sgn}\left( \xi_{i_F}^{\bar{x},2} + \mathbf{2}\,\xi_{i_F}^{x,2} + \xi_{i_F}^{x,1} + \sum_{\mu=3}^{P} \xi_{i_F}^{x,\mu} \right)
    \end{align}
    \textbf{The ``Broken Symmetry'':} $\xi_{i_F}^{x,2}$ is more likely; it will influence $S_{i_F}^F(1)$, \textbf{but} produces a very disturbed output (\textit{too disturbed $\implies$ very light disambiguation}).
\end{enumerate}




% !TEX encoding = UTF-8
% !TEX program = pdflatex
% !TEX root = ../A-thesis-main.tex

%\lipsum
\section{DMFT}
I would like to reuse the structure/framework for the analytical path-integral approach to DMFT. \cite{metz_dynamical_2025}



\chapter{Numerical Simulations}
% !TEX encoding = UTF-8
% !TEX program = pdflatex
% !TEX root = ../A-thesis-main.tex








\chapter{Results}
% !TEX encoding = UTF-8
% !TEX program = pdflatex
% !TEX root = ../A-thesis-main.tex

In conclusione abbiamo imparato cose belle.\\
Vediamo se funziona fare una citazione \parencite{pippo}.
\lipsum


\appendix
%\chapter{Ito calculus}
% !TEX encoding = UTF-8
% !TEX program = pdflatex
% !TEX root = ../A-thesis-main.tex

qualcosa da aggiungere?\\
\lipsum

\printbibliography[heading=bibintoc, title={References}]


\end{document}